\chapter{Solvable Groups and \texorpdfstring{$p$}{p}-Groups}

\section{Derived Series and Solvable Groups}

\begin{definition}[Commutator subgroup]
Let \(G\) be a group. The \textbf{commutator subgroup} of \(G\) is
\[
    G'=[G,G]
    =
    \gen{[x,y]:x,y\in G},
\]
where
\[
    [x,y]=x^{-1}y^{-1}xy.
\]
Equivalently, \(G'\) is the smallest normal subgroup \(N\lhd G\) such that
\[
    G/N
\]
is abelian.
\end{definition}

\begin{definition}[Derived series]
Define the derived series of \(G\) by
\[
    G^{(0)}=G,
    \qquad
    G^{(1)}=G',
\]
and inductively
\[
    G^{(n)}=\bigl(G^{(n-1)}\bigr)'
    \qquad(n\geq 1).
\]
Thus we have a descending normal series
\[
    G\trianglerighteq G'\trianglerighteq G''\trianglerighteq
    \cdots
    \trianglerighteq G^{(n)}
    \trianglerighteq \cdots .
\]
\end{definition}

\begin{example}[Some commutator subgroups]
\begin{enumerate}[label=(\arabic*)]
    \item If \(G\) is simple and nonabelian, then
    \[
        G'=G.
    \]
    Indeed, \(G'\lhd G\). Since \(G\) is nonabelian, \(G'\neq 1\), and hence by simplicity
    \(G'=G\).

    \item For the symmetric group \(S_n\), one has
    \[
        S_n'=A_n
        \qquad(n\geq 3).
    \]
    In particular,
    \[
        S_4'=A_4,
        \qquad
        S_n'=A_n \quad(n\geq 5).
    \]
    This is because \(A_n\) is generated by \(3\)-cycles, and \(3\)-cycles occur as commutators.
\end{enumerate}
\end{example}

\begin{remark}
If \(G\) is finite, then the descending chain
\[
    G\trianglerighteq G'\trianglerighteq G^{(2)}
    \trianglerighteq \cdots
\]
must eventually stabilize. Hence either
\[
    G^{(n)}=1
\]
for some \(n\), or else
\[
    G^{(n)}=G^{(n+1)}=G^{(n+2)}=\cdots
\]
for all sufficiently large \(n\). The stable term is sometimes denoted by
\[
    G^{(\infty)}.
\]
\end{remark}

\begin{definition}[Solvable and perfect groups]
Let \(G\) be a group.
\begin{enumerate}[label=(\arabic*)]
    \item We say that \(G\) is \textbf{solvable} if
    \[
        G^{(n)}=1
    \]
    for some \(n\geq 0\).

    \item We say that \(G\) is \textbf{perfect} if
    \[
        G=G'.
    \]
\end{enumerate}
\end{definition}

\begin{example}[Perfect groups]
The following are important examples of perfect groups.
\begin{enumerate}[label=(\arabic*)]
    \item Every nonabelian simple group is perfect.

    \item The alternating group \(A_n\) is perfect for \(n\geq 5\). For example,
    \[
        (123)(345)(123)^{-1}(345)^{-1}=(143).
    \]
    Since \(A_n\) is generated by \(3\)-cycles, this implies
    \[
        A_n'=A_n.
    \]

    \item The group \(\operatorname{SL}_d(q)\) is perfect whenever it is nonsolvable; more precisely,
    \[
        \operatorname{SL}_d(q)'=\operatorname{SL}_d(q)
    \]
    for \(d\geq 2\), except in the small cases
    \[
        (d,q)=(2,2),(2,3).
    \]
    Its center is
    \[
        Z(\operatorname{SL}_d(q))
        =
        \{\lambda I_d:\lambda^d=1,\ \lambda\in \mathbb F_q^\times\},
    \]
    and hence
    \[
        |Z(\operatorname{SL}_d(q))|
        =
        \gcd(d,q-1).
    \]
\end{enumerate}
\end{example}

\section{Basic Properties of Solvable Groups}

\begin{proposition}[Properties of solvable groups]
Let \(G\) be a group.
\begin{enumerate}[label=(\arabic*)]
    \item If \(G\) is solvable, then every subgroup of \(G\) and every quotient group of \(G\) is solvable.

    \item If \(N\lhd G\), and both \(N\) and \(G/N\) are solvable, then \(G\) is solvable.

    \item Every finite \(p\)-group is solvable.
\end{enumerate}
\end{proposition}

\begin{proof}
\begin{enumerate}[label=(\arabic*)]
    \item Let \(H\leq G\). Since
    \[
        H'\leq G',
    \]
    induction gives
    \[
        H^{(n)}\leq G^{(n)}
    \]
    for all \(n\). Thus if \(G^{(n)}=1\), then \(H^{(n)}=1\), so \(H\) is solvable.

    Now let \(N\lhd G\). The natural map \(G\to G/N\) sends \(G^{(i)}\) onto a subgroup of
    \((G/N)^{(i)}\), and in fact
    \[
        (G/N)^{(i)}
        =
        G^{(i)}N/N.
    \]
    Hence if \(G^{(n)}=1\), then
    \[
        (G/N)^{(n)}=1,
    \]
    so \(G/N\) is solvable.

    \item Since \(G/N\) is solvable, there exists \(r\) such that
    \[
        (G/N)^{(r)}=1.
    \]
    Hence
    \[
        G^{(r)}\leq N.
    \]
    Since \(N\) is solvable, there exists \(s\) such that
    \[
        N^{(s)}=1.
    \]
    Therefore
    \[
        G^{(r+s)}
        \leq N^{(s)}
        =
        1.
    \]
    Thus \(G\) is solvable.

    \item We argue by induction on \(|G|\). If \(G=1\), there is nothing to prove.

    If \(G\neq 1\) is a finite \(p\)-group, then
    \[
        Z(G)\neq 1.
    \]
    Since \(Z(G)\lhd G\) is abelian, it is solvable. Also \(G/Z(G)\) is a smaller \(p\)-group, hence solvable by induction. By part (2), \(G\) is solvable.
\end{enumerate}
\end{proof}

\section{Double Cosets and a Normalizer Lemma}

\begin{definition}[Double coset]
Let \(P\leq G\). For \(x\in G\), the subset
\[
    PxP
    =
    \{h_1xh_2:h_1,h_2\in P\}
\]
is called a \textbf{double coset} of \(P\) in \(G\).
\end{definition}

\begin{lemma}[Right cosets inside a double coset]
Let \(P\leq G\), and write
\[
    P^x=x^{-1}Px.
\]
Then the number of right cosets of \(P\) contained in the double coset \(PxP\) is
\[
    \frac{|P|}{|P^x\cap P|}
    =
    [P:P^x\cap P].
\]
\end{lemma}

\begin{proof}
We may write
\[
    PxP=\bigcup_{h\in P}Pxh.
\]
Two right cosets \(Pxh_1\) and \(Pxh_2\) are equal if and only if
\[
    Pxh_1=Pxh_2.
\]
This is equivalent to
\[
    xh_1h_2^{-1}x^{-1}\in P,
\]
or equivalently
\[
    h_1h_2^{-1}\in x^{-1}Px=P^x.
\]
Since \(h_1h_2^{-1}\in P\), this condition becomes
\[
    h_1h_2^{-1}\in P^x\cap P.
\]
Therefore the distinct right cosets \(Pxh\) are parametrized by the right cosets of
\(P^x\cap P\) in \(P\). Hence the number is
\[
    [P:P^x\cap P]
    =
    \frac{|P|}{|P^x\cap P|}.
\]
\end{proof}

\begin{lemma}[Normalizer lemma]
Let \(G\) be finite, and let \(P\) be a \(p\)-subgroup of \(G\). If \(P\) is not a Sylow \(p\)-subgroup of \(G\), then
\[
    P<N_G(P).
\]
\end{lemma}

\begin{proof}
Decompose \(G\) into double cosets of \(P\):
\[
    G
    =
    \bigsqcup_{i=0}^{t}Px_iP,
\]
where we take
\[
    x_0=1,
    \qquad
    Px_0P=P.
\]
Thus
\[
    P
\]
is the unique right coset of \(P\) inside \(P\).

Dividing the double coset decomposition by right cosets of \(P\), we obtain
\[
    [G:P]
    =
    1+
    \sum_{i=1}^{t}
    [P:P^{x_i}\cap P].
\]
Since \(P\) is not Sylow, the index \([G:P]\) is divisible by \(p\).

Each summand
\[
    [P:P^{x_i}\cap P]
\]
is a power of \(p\). Hence it is either \(1\), or divisible by \(p\). If all of the summands with \(i\geq 1\) were divisible by \(p\), then the right-hand side would be congruent to \(1\pmod p\), contradicting \(p\mid [G:P]\). Therefore for some \(i\geq 1\),
\[
    [P:P^{x_i}\cap P]=1.
\]
Thus
\[
    P^{x_i}\cap P=P,
\]
so
\[
    P^{x_i}=P.
\]
Therefore
\[
    x_i\in N_G(P).
\]
Since \(i\geq 1\), we may choose \(x_i\notin P\). Hence
\[
    P<N_G(P).
\]
\end{proof}

\section{Subgroups of \texorpdfstring{$p$}{p}-Groups}

\begin{theorem}[Subgroups and normal subgroups of \(p\)-groups]
Let \(G\) be a finite \(p\)-group.
\begin{enumerate}[label=(\arabic*)]
    \item If \(H<G\), then
    \[
        H<N_G(H).
    \]

    \item If \(H\) is a maximal subgroup of \(G\), then
    \[
        H\lhd G
        \qquad\text{and}\qquad
        |G:H|=p.
    \]

    \item The group \(G\) is solvable.

    \item The center of \(G\) is nontrivial:
    \[
        Z(G)>1.
    \]

    \item If \(N\lhd G\) and \(|N|=p\), then
    \[
        N\leq Z(G).
    \]
\end{enumerate}
\end{theorem}

\begin{proof}
\begin{enumerate}[label=(\arabic*)]
    \item Since \(G\) itself is a \(p\)-group, every proper subgroup \(H<G\) is not a Sylow \(p\)-subgroup of \(G\). By the normalizer lemma,
    \[
        H<N_G(H).
    \]

    \item If \(H\) is maximal, then by part (1),
    \[
        H<N_G(H)\leq G.
    \]
    Hence maximality forces
    \[
        N_G(H)=G.
    \]
    Thus \(H\lhd G\). The quotient \(G/H\) has no nontrivial proper subgroups, because \(H\) is maximal. Since \(G/H\) is a nontrivial \(p\)-group, it must have order \(p\). Hence
    \[
        |G:H|=p.
    \]

    \item This follows by induction on \(|G|\). Choose a maximal subgroup \(H<G\). By part (2),
    \[
        H\lhd G
        \qquad\text{and}\qquad
        G/H\cong C_p.
    \]
    By induction, \(H\) is solvable, and \(G/H\) is abelian, hence solvable. Therefore \(G\) is solvable.

    \item Let \(G\) act on itself by conjugation. The class equation is
    \[
        |G|
        =
        |Z(G)|
        +
        \sum_i [G:C_G(x_i)],
    \]
    where the \(x_i\)'s represent the noncentral conjugacy classes.

    If \(x_i\notin Z(G)\), then
    \[
        C_G(x_i)<G,
    \]
    so
    \[
        [G:C_G(x_i)]
    \]
    is divisible by \(p\). Since \(|G|\) is also divisible by \(p\), it follows that
    \[
        p\mid |Z(G)|.
    \]
    Thus \(Z(G)\neq 1\).

    \item Since \(N\lhd G\), the normalizer-centralizer lemma gives
    \[
        G/C_G(N)
        \hookrightarrow
        \Aut(N).
    \]
    But \(|N|=p\), so
    \[
        N\cong C_p
        \qquad\text{and}\qquad
        \Aut(N)\cong C_{p-1}.
    \]
    Hence
    \[
        |G:C_G(N)|
        \mid p-1.
    \]
    On the other hand, \(G/C_G(N)\) is a quotient of the \(p\)-group \(G\), so its order is a power of \(p\). Therefore
    \[
        |G:C_G(N)|=1.
    \]
    Thus
    \[
        G=C_G(N),
    \]
    which means every element of \(G\) centralizes \(N\). Hence
    \[
        N\leq Z(G).
    \]
\end{enumerate}
\end{proof}

\begin{lemma}[Cyclic quotient by the center]
    Let \(G\) be a group. If
    \[
        G/Z(G)
    \]
    is cyclic, then \(G\) is abelian.
    \end{lemma}
    
    \begin{proof}
    Assume
    \[
        G/Z(G)=\gen{gZ(G)}
    \]
    for some \(g\in G\). Then every element of \(G\) can be written in the form
    \[
        g^m z
    \]
    for some \(m\in \mathbb Z\) and some \(z\in Z(G)\).
    
    Let \(x,y\in G\). Then there exist \(m,n\in\mathbb Z\) and
    \(z_1,z_2\in Z(G)\) such that
    \[
        x=g^m z_1,
        \qquad
        y=g^n z_2.
    \]
    Since \(z_1,z_2\in Z(G)\), they commute with all elements of \(G\). Hence
    \[
        xy
        =
        (g^m z_1)(g^n z_2)
        =
        g^{m+n}z_1z_2
        =
        g^{n+m}z_2z_1
        =
        (g^n z_2)(g^m z_1)
        =
        yx.
    \]
    Thus any two elements of \(G\) commute, so \(G\) is abelian.
    \end{proof}

\begin{corollary}
Every group of order \(p^2\), where \(p\) is prime, is abelian.
\end{corollary}

\begin{proof}
Let \(|G|=p^2\). Since \(G\) is a \(p\)-group,
\[
    Z(G)\neq 1.
\]
Thus
\[
    |Z(G)|=p
    \qquad\text{or}\qquad
    |Z(G)|=p^2.
\]
If \(|Z(G)|=p^2\), then \(Z(G)=G\), so \(G\) is abelian.

If \(|Z(G)|=p\), then
\[
    |G/Z(G)|=p,
\]
so \(G/Z(G)\) is cyclic. But if \(G/Z(G)\) is cyclic, then \(G\) is abelian. Hence \(Z(G)=G\), contradiction. Therefore \(G\) must be abelian.
\end{proof}

\section{Groups of Order \texorpdfstring{$p^3$}{p3}}

\begin{theorem}[Structure of groups of order \(p^3\)]
Let \(G\) be a group of order \(p^3\), where \(p\) is prime. Then either \(G\) is abelian, in which case
\[
    G\cong C_{p^3},
    \qquad
    G\cong C_{p^2}\times C_p,
    \qquad\text{or}\qquad
    G\cong C_p\times C_p\times C_p,
\]
or \(G\) is nonabelian. In the nonabelian case, the following possibilities occur.

\begin{enumerate}[label=(\arabic*)]
    \item If \(p=2\), then
    \[
        G\cong D_8
        \qquad\text{or}\qquad
        G\cong Q_8.
    \]

    \item If \(p\) is odd, then either
    \[
        G\cong
        \gen{
            a,b
            \mid
            a^{p^2}=b^p=1,\ a^b=a^{1+p}
        },
    \]
    or
    \[
        G\cong
        \gen{
            a,b,c
            \mid
            a^p=b^p=c^p=1,\ c=[a,b],\ [a,c]=[b,c]=1
        }.
    \]
    The first group has exponent \(p^2\), and the second group has exponent \(p\).
\end{enumerate}
\end{theorem}

\begin{proof}
Assume \(G\) is nonabelian.

Since \(G\) is a \(p\)-group, \(Z(G)\neq 1\). If \(|Z(G)|=p^2\), then
\[
    |G/Z(G)|=p,
\]
so \(G/Z(G)\) is cyclic, forcing \(G\) to be abelian. This is impossible. Hence
\[
    |Z(G)|=p.
\]
Also, \(G/Z(G)\) has order \(p^2\), so it is abelian. Therefore
\[
    G'\leq Z(G).
\]
Since \(G\) is nonabelian, \(G'\neq 1\), and hence
\[
    G'=Z(G)\cong C_p.
\]

We split into two cases.

\medskip

\noindent\textbf{Case 1: \(G\) contains an element of order \(p^2\).}

Let \(a\in G\) with
\[
    |a|=p^2.
\]
Then
\[
    \gen{a}
\]
has order \(p^2\), so it is a maximal subgroup of \(G\). Hence
\[
    \gen{a}\lhd G.
\]
Choose
\[
    b\in G\setminus \gen{a}.
\]
Then \(G/\gen{a}\) has order \(p\), so
\[
    b^p\in \gen{a}.
\]

Conjugation by \(b\) gives an automorphism of \(\gen{a}\). Thus
\[
    a^b=a^\lambda
\]
for some integer \(\lambda\). Since \(b^p\in\gen{a}\), the induced automorphism has order dividing \(p\), so
\[
    \lambda^p\equiv 1\pmod{p^2}.
\]

If \(p=2\), then the only nontrivial automorphism of \(C_4\) is inversion. Hence
\[
    a^b=a^{-1}.
\]
If there exists \(b\notin\gen{a}\) with \(b^2=1\), then
\[
    G\cong D_8.
\]
Otherwise every element outside \(\gen{a}\) has order \(4\), and one obtains
\[
    G\cong Q_8.
\]

Now assume \(p\) is odd. Since \(G\) is nonabelian, the action is nontrivial, so
\[
    \lambda\not\equiv 1\pmod{p^2}.
\]
From
\[
    \lambda^p\equiv 1\pmod{p^2},
\]
we get
\[
    \lambda\equiv 1\pmod p.
\]
Thus
\[
    \lambda=1+kp
\]
for some \(k\) with \(1\leq k\leq p-1\).

Write
\[
    b^p=a^{\ell p}
\]
for some integer \(\ell\). A direct computation gives
\[
    (ba^{-\ell})^p=1.
\]
Replacing \(b\) by \(ba^{-\ell}\), we may assume
\[
    b^p=1.
\]
Finally, replacing \(b\) by a suitable power of \(b\), we may arrange that
\[
    a^b=a^{1+p}.
\]
Therefore
\[
    G\cong
    \gen{
        a,b
        \mid
        a^{p^2}=b^p=1,\ a^b=a^{1+p}
    }.
\]

\medskip

\noindent\textbf{Case 2: \(G\) has no element of order \(p^2\).}

Then every nonidentity element of \(G\) has order \(p\). If \(p=2\), then every element has order \(2\), and hence \(G\) is abelian, a contradiction. Therefore \(p\) is odd.

Choose \(z\in Z(G)\) with \(z\neq 1\). Then
\[
    |\gen{z}|=p,
\]
and
\[
    |G/\gen{z}|=p^2.
\]
Hence \(G/\gen{z}\) is abelian, so
\[
    G'\leq \gen{z}.
\]
Since \(G\) is nonabelian, \(G'\neq 1\). Thus
\[
    G'=\gen{z}=Z(G).
\]

Choose \(a,b\in G\) such that
\[
    [a,b]\neq 1.
\]
Let
\[
    c=[a,b].
\]
Then
\[
    c\in G'=Z(G),
\]
so \(c\) is central. Since every nonidentity element has order \(p\),
\[
    a^p=b^p=c^p=1.
\]
Moreover, \(\gen{a,b}\) is not abelian, so it cannot have order \(p\) or \(p^2\). Hence
\[
    \gen{a,b}=G.
\]
Thus
\[
    G\cong
    \gen{
        a,b,c
        \mid
        a^p=b^p=c^p=1,\ c=[a,b],\ [a,c]=[b,c]=1
    }.
\]
This is the nonabelian group of order \(p^3\) and exponent \(p\).
\end{proof}

\section{Automorphisms of Nonabelian Groups of Order \texorpdfstring{$p^3$}{p3}}

\begin{remark}
For the two nonabelian groups of order \(8\), one has
\[
    \Aut(D_8)\cong D_8,
\]
and
\[
    \Aut(Q_8)\cong S_4.
\]
\end{remark}

\begin{definition}[The groups \(P_-^{1+2}\) and \(P_+^{1+2}\)]
For odd \(p\), denote the two nonabelian groups of order \(p^3\) by
\[
    P_-^{1+2}
    =
    \gen{
        a,b
        \mid
        a^{p^2}=b^p=1,\ a^b=a^{1+p}
    },
\]
and
\[
    P_+^{1+2}
    =
    \gen{
        a,b,c
        \mid
        a^p=b^p=c^p=1,\ c=[a,b],\ [a,c]=[b,c]=1
    }.
\]
Then in both cases
\[
    G'=Z(G)\cong C_p.
\]
\end{definition}

\begin{remark}[Induced automorphisms on \(G/G'\)]
Let \(G\) be one of the two nonabelian groups of order \(p^3\), with \(p\) odd. Since
\[
    G'=Z(G)\cong C_p,
\]
every automorphism \(\alpha\in\Aut(G)\) satisfies
\[
    (G')^\alpha=G',
    \qquad
    Z(G)^\alpha=Z(G).
\]
Hence \(\alpha\) induces a linear automorphism
\[
    \overline{\alpha}\in \Aut(G/G')
    \cong \operatorname{GL}(2,p).
\]

For
\[
    G=P_-^{1+2},
\]
we have
\[
    G=\gen{a}:\gen{b},
    \qquad
    G'=\gen{a^p}.
\]
Indeed, from
\[
    a^b=a^{1+p},
\]
we get
\[
    (a^p)^b=(a^b)^p=(a^{1+p})^p=a^p.
\]
The automorphism group has the form
\[
    \Aut(P_-^{1+2})\cong P:C_{p-1},
\]
where \(P\) is a \(p\)-group.

For
\[
    G=P_+^{1+2},
\]
the commutator map induces a nondegenerate alternating bilinear form
\[
    (G/G')\times (G/G')\longrightarrow G',
    \qquad
    (\overline{x},\overline{y})\longmapsto [x,y].
\]
Automorphisms fixing \(G'\) pointwise induce elements of
\[
    \operatorname{Sp}(2,p)\cong \operatorname{SL}(2,p).
\]
Allowing scalar multiplication on \(G'\) extends this to
\[
    \operatorname{GL}(2,p).
\]
Equivalently,
\[
    \Aut(P_+^{1+2})/O_p(\Aut(P_+^{1+2}))
    \cong
    \operatorname{GL}(2,p).
\]
More explicitly,
\[
    \Aut(P_+^{1+2})
    \cong
    C_p^2:\operatorname{GL}(2,p).
\]
\end{remark}